\usepackage[top=3cm,bottom=3cm,left=3cm,right=3cm,headsep=10pt,a4paper]{geometry} % Page margins

\usepackage{graphicx} % Required for including pictures

\usepackage{lipsum} % Inserts dummy text

\usepackage{tikz}                                                            % Required for drawing custom shapes

\usepackage[italian]{babel}

\usepackage{enumitem} % Customize lists
\setlist{nolistsep} % Reduce spacing between bullet points and numbered lists

\usepackage{booktabs} % Required for nicer horizontal rules in tables

\usepackage{xcolor} % Required for specifying colors by name
\definecolor{mainColor}{RGB}{243,102,25} % Define the orange color used for highlighting throughout the book

\usepackage{hyperref}

\usepackage{avant} % Use the Avantgarde font for headings
%\usepackage{times} % Use the Times font for headings
\usepackage{mathptmx} % Use the Adobe Times Roman as the default text font together with math symbols from the Sym­bol, Chancery and Com­puter Modern fonts

\usepackage{microtype} % Slightly tweak font spacing for aesthetics
\usepackage[utf8]{inputenc} % Required for including letters with accents
\usepackage[T1]{fontenc} % Use 8-bit encoding that has 256 glyphs

%----------------------------------------------------------------------------------------
%	BIBLIOGRAPHY AND INDEX
%----------------------------------------------------------------------------------------

% \usepackage{csquotes}
% \usepackage[style=alphabetic,citestyle=numeric,sorting=nyt,sortcites=true,autopunct=true,autolang=hyphen,hyperref=true,abbreviate=false,backref=true,backend=biber,defernumbers=true]{biblatex}
% \addbibresource{bibliography.bib} % BibTeX bibliography file
% \defbibheading{bibempty}{}

\usepackage{calc} % For simpler calculation - used for spacing the index letter headings correctly
% \usepackage{makeidx} % Required to make an index
% \makeindex % Tells LaTeX to create the files required for indexing


%----------------------------------------------------------------------------------------
%	MAIN TABLE OF CONTENTS
%----------------------------------------------------------------------------------------

\usepackage{titletoc}                                                        % Required for manipulating the table of contents

\contentsmargin{0cm}                                                         % Removes the default margin

% Part text styling
\titlecontents{part}[1cm]                                                    % Indentation
{\addvspace{20pt}\centering\LARGE\bfseries}                                  % Spacing and font options for parts
{}
{}  
{}

% Chapter text styling
\titlecontents{chapter}[1.25cm]                                              % Indentation
{\addvspace{12pt}\large\sffamily\bfseries}                                   % Spacing and font options for chapters
{\color{mainColor!60}\contentslabel[\Large\thecontentslabel]{1.25cm}\color{mainColor}} % Chapter number
{\color{mainColor}}  
{\color{mainColor!60}\normalsize\;\titlerule*[.5pc]{.}\;\thecontentspage}         % Page number

% Section text styling
\titlecontents{section}[1.25cm]                                              % Indentation
{\addvspace{3pt}\sffamily\bfseries}                                          % Spacing and font options for sections
{\contentslabel[\thecontentslabel]{1.25cm}}                                  % Section number
{}
{\hfill\color{black}\thecontentspage}                                        % Page number
[]

% Subsection text styling
\titlecontents{subsection}[1.25cm]                                           % Indentation
{\addvspace{1pt}\sffamily\small}                                             % Spacing and font options for subsections
{\contentslabel[\thecontentslabel]{1.25cm}}                                  % Subsection number
{}
{\ \titlerule*[.5pc]{.}\;\thecontentspage}                                   % Page number
[]

%----------------------------------------------------------------------------------------
%	MINI TABLE OF CONTENTS IN PART HEADS
%----------------------------------------------------------------------------------------

% Chapter text styling
\titlecontents{lchapter}[0em]                                                % Indenting
{\addvspace{15pt}\large\sffamily\bfseries}                                   % Spacing and font options for chapters
{\color{mainColor}\contentslabel[\Large\thecontentslabel]{1.25cm}\color{mainColor}} % Chapter number
{}  
{\color{mainColor}\normalsize\sffamily\bfseries\;\titlerule*[.5pc]{.}\;\thecontentspage} % Page number

% Section text styling
\titlecontents{lsection}[0em]                                                % Indenting
{\sffamily\small}                                                            % Spacing and font options for sections
{\contentslabel[\thecontentslabel]{1.25cm}}                                  % Section number
{}
{}


%----------------------------------------------------------------------------------------
%	PAGE HEADERS
%----------------------------------------------------------------------------------------

\usepackage{fancyhdr} % Required for header and footer configuration

\pagestyle{fancy}
\renewcommand{\chaptermark}[1]{\markboth{\sffamily\normalsize\bfseries\chaptername\ \thechapter.\ #1}{}} 
\renewcommand{\sectionmark}[1]{\markright{\sffamily\normalsize\thesection\hspace{5pt}#1}{}}

\fancyhf{} \fancyhead[LE,RO]{\sffamily\normalsize\thepage}         % Font setting for the page number in the header
\fancyhead[LO]{\rightmark}                                         % Print the nearest section name on the left side of odd pages
\fancyhead[RE]{\leftmark}                                          % Print the current chapter name on the right side of even pages

\renewcommand{\headrulewidth}{0.5pt}                               % Width of the rule under the header
\addtolength{\headheight}{2.5pt}                                   % Increase the spacing around the header slightly
\renewcommand{\footrulewidth}{0pt}                                 % Removes the rule in the footer

\fancypagestyle{plain}{\fancyhead{}\renewcommand{\headrulewidth}{0pt}} % Style for when a plain pagestyle is specified

% Removes the header from odd empty pages at the end of chapters
\makeatletter
\renewcommand{\cleardoublepage}{
    \clearpage
    \ifodd
        \c@page
    \else
        \hbox{}
        \vspace*{\fill}
        \thispagestyle{empty}
        \newpage
    \fi
}

%----------------------------------------------------------------------------------------
%	PART HEADINGS
%----------------------------------------------------------------------------------------

\newlength\esp
\setlength\esp{4pt}
\newlength\ecart
\setlength\ecart{1.2cm-\esp}
\newcommand{\thepartimage}{}
\newcommand{\partimage}[1]{\renewcommand{\thepartimage}{#1}}
\def\@part[#1]#2{
    \ifnum \c@secnumdepth >-2\relax
        \refstepcounter{part}
    \fi
    \startcontents
    \markboth{}{}
    {
        \thispagestyle{empty}
        \begin{tikzpicture}[remember picture,overlay]
            \node at (current page.north west){
                \begin{tikzpicture}[remember picture,overlay]
                    \fill[mainColor!10](0cm,0cm) rectangle (\paperwidth,-\paperheight);             % Sfondo PART
                    
                    \ifx\thepartimage\@empty
                    \else
                        \node[anchor=north west, inner sep=0pt,opacity=.75] at (0,0) {              % Sfondo PART (partimage)         
                            \includegraphics[width=\paperwidth]{\thepartimage}
                        };  
                    \fi

                    \node[anchor=north west] at (1cm,-2.75cm){                                      % Numero PART
                        \color{white!80}\fontsize{250}{100}\sffamily\bfseries\@Roman\c@part
                    }; 
                    \node[anchor=north east] at (\paperwidth-1.5cm,-11cm){                          % Titolo PART
                        \parbox[t][][t]{\paperwidth-3cm}{
                            \strut\centering\color{mainColor!70}\fontsize{50}{50}\sffamily\bfseries#2
                        }
                    };
                    \node[anchor=north east] at (\paperwidth-1cm,-\paperheight+15cm){               % Indice PART        
                        \parbox[t][][t]{8.5cm}{
                            \printcontents{l}{0}{\setcounter{tocdepth}{1}}
                        }
                    };
                \end{tikzpicture}
            };
        \end{tikzpicture}
    }
    \@endpart
}

\def\@endpart{
    \vfil
    \newpage
    \if@twoside
        \if@openright
            \null
            \thispagestyle{empty}
            \newpage
        \fi
    \fi
    \if@tempswa
        \twocolumn
    \fi
}

%----------------------------------------------------------------------------------------
%	CHAPTER HEADINGS
%----------------------------------------------------------------------------------------

\newcommand{\thechapterimage}{}
\newcommand{\chapterimage}[1]{\renewcommand{\thechapterimage}{#1}}
\def\@makechapterhead#1{                            % Capitoli numerati
    {
        \parindent \z@ \raggedright \normalfont
        \begin{tikzpicture}[remember picture,overlay]
            \node at (current page.north west){
                \begin{tikzpicture}[remember picture,overlay]
                    \node[anchor=north west,inner sep=0pt] at (0,0) {                           % Sfondo capitolo (img)
                        \ifx\thechapterimage\@empty
                        \else
                            \includegraphics[width=\paperwidth]{\thechapterimage}
                        \fi
                    };
                    \draw[anchor=west] (\Gm@lmargin,-9cm) node [line width=2pt,rounded corners=15pt,draw=mainColor,fill=white,fill opacity=0.5,inner sep=15pt]{
                        \strut\makebox[22cm]{}
                    };
                    \node[anchor=north west] at (2cm,-2.75cm){                                  % Numero capitolo
                        \color{mainColor!40}\fontsize{150}{100}\sffamily\bfseries\thechapter
                    }; 
                    \draw[anchor=west] (\Gm@lmargin+.3cm,-9cm) node {
                        \huge\sffamily\bfseries\color{black}#1\strut
                    };
                \end{tikzpicture}
            };
        \end{tikzpicture}
        \par\vspace*{270\p@}
    }
}

\def\@makeschapterhead#1{%          % Capitoli non numerati
    \begin{tikzpicture}[remember picture,overlay]
        \node at (current page.north west){
            \begin{tikzpicture}[remember picture,overlay]
                \node[anchor=north west,inner sep=0pt] at (0,0) {
                    \ifx\thechapterimage\@empty
                    \else
                        \includegraphics[width=\paperwidth]{\thechapterimage}
                    \fi
                };
                \draw[anchor=west] (\Gm@lmargin,-9cm) node [line width=2pt,rounded corners=15pt,draw=mainColor,fill=white,fill opacity=0.5,inner sep=15pt]{
                    \strut\makebox[22cm]{}
                };
                \draw[anchor=west] (\Gm@lmargin+.3cm,-9cm) node {
                    \huge\sffamily\bfseries\color{black}#1\strut
                };
            \end{tikzpicture}
        };
    \end{tikzpicture}
    \par\vspace*{270\p@}
}
\makeatother

%----------------------------------------------------------------------------------------
%	HYPERLINKS IN THE DOCUMENTS
%----------------------------------------------------------------------------------------

\hypersetup{hidelinks,colorlinks=false,breaklinks=true,urlcolor= mainColor,bookmarksopen=false,pdftitle={Title},pdfauthor={Author}}
\usepackage{bookmark}
\bookmarksetup{
    open,
    numbered,
    addtohook={%
        \ifnum\bookmarkget{level}=0 % chapter
            \bookmarksetup{bold}%
        \fi
        \ifnum\bookmarkget{level}=-1 % part
            \bookmarksetup{color=mainColor,bold}%
        \fi
    }
}


%----------------------------------------------------------------------------------------
%	ESERCIZI
%----------------------------------------------------------------------------------------

\usepackage[most]{tcolorbox}
\usepackage{xcolor}
\usepackage{tikz}
\usepackage{ifthen}

\newcounter{esercizio}[section]                                    % Contatore progressivo per sezione
\renewcommand{\theesercizio}{\arabic{esercizio}}

\newenvironment{esercizio}[1][]{
    \vspace{0.25cm}
    \refstepcounter{esercizio}
    \noindent
    \begin{minipage}[t]{1.8cm}
        \hspace{0.1cm}
        \begin{tikzpicture}
            \node[
                fill=mainColor!60, 
                text=white, 
                font=\bfseries, 
                inner sep=0,
                minimum width=1.2cm, 
                minimum height=0.6cm,
                rounded corners=0.1cm
            ] (n) at (0,0) {\theesercizio};
            \foreach \i in {1,...,#1} {
                \node[
                    circle, 
                    fill=mainColor!90,
                    inner sep=0, 
                    minimum size=0.2cm
                ] at (0.3*\i-0.6, -0.45) {};
            }
        \end{tikzpicture}\par
    \end{minipage}
    \begin{minipage}[t]{\dimexpr\linewidth-2cm\relax}
    \vspace{-.8cm}
}{
    \end{minipage}
    \par
}